\begin{description}
\item[(a)] Taking multiple reference points, we get that the BLR emission
  lags by about \textbf{20--25 days}. \hfill(1.0 points)

  Answers from 15 to 25 days are acceptable.

\item[(b)] As the AGN is located very far, the time lag can be approximated
  as the time taken by UV emission to reach BLR region. Thus,
  \begin{align*}
    R_{\mathrm{BLR}} = c\Delta t
      &= 3\times 10^{8}\times 20\times 86\,400 \\
      &= 5.2\times 10^{14}\,\mathrm{m} = (0.017\pm 0.004)\,\mathrm{pc}
  \end{align*}
  \hfill(3.0 points)

\item[(c)] As this AGN is 100\,Mpc away from us,
  \begin{align*}
    \theta_{\mathrm{BLR}} &= \frac{0.017}{100e6}\times 206265 \\
      &= (3.5\pm 0.9)\times 10^{-5}\,\mathrm{arcsec}
  \end{align*}
  \hfill(2.0 points)

\item[(d)] The FWHM is approximately $(85\pm 5)$\,\AA\ and the peak is
  approximately at $(4940\pm 5)$\,\AA. \hfill(2.0 points)
  \[ v_\sigma = \frac{\sigma c}{\lambda_{\mathrm{peak}}}
     = \frac{\mathrm{FWHM}\times c}{2.35\,\lambda_{\mathrm{peak}}}
     = \frac{85\times 3\times 10^{5}}{2.35\times 4940} \]
  \hfill(2.0 points)
  \[ v_\sigma = (2200\pm 140)\,\mathrm{km\,s^{-1}} \]
  \hfill(1.0 points)

\item[(e)] \mbox{}
  \begin{align*}
    M_{\mathrm{vir,BH}} &= \frac{v_\sigma^2 R_{\mathrm{BLR}}}{G}
      = \frac{(2.2\times 10^{6})^2\times 5.2\times 10^{14}}
             {6.674\times 10^{-11}} \\
      &= 3.8\times 10^{37}\,\mathrm{kg} \\
    M_{\mathrm{vir,BH}} &= (1.9\pm 0.7)\times 10^{7}\,M_\odot
  \end{align*}
  \hfill(4.0 points)
\end{description}
